\documentclass[11pt,a4paper]{article}
\usepackage[utf8]{inputenc}
\usepackage[T1]{fontenc}
\usepackage{lmodern}
\usepackage[german]{babel}
\usepackage{amsmath,amssymb,amsthm,amsfonts}
\usepackage{geometry}
\usepackage{microtype}
\usepackage{hyperref}
\usepackage{cite}

\geometry{margin=1in}

% Theorem-Umgebungen auf Deutsch
\newtheorem{theorem}{Satz}[section]
\newtheorem{lemma}[theorem]{Lemma}
\newtheorem{definition}[theorem]{Definition}
\newtheorem{proposition}[theorem]{Proposition}
\newtheorem{corollary}[theorem]{Korollar}
\theoremstyle{definition}
\newtheorem{remark}[theorem]{Bemerkung}

\title{\textbf{Die arithmetische Starrheit des Collatz-Graphen:}\\[0.5em]
	\Large Eine quaternionisch-primalgeometrische und adisch-analytische Fundamentalerklärung der Syracuse-Iteration}

\author{\textbf{Thomas Hoffbauer}\thanks{Bamberg, Deutschland. E-Mail: \texttt{thomas.hoffbauer@bamberger-modell.de}}}
\date{September 2026}

\begin{document}
	
	\maketitle
	
	\begin{abstract}
		Die Collatz-Vermutung (Syracuse-Problem) widersteht seit fast einem Jahrhundert einer vollständigen Auflösung, primär weil konventionelle dynamische und stochastische Methoden die Vorwärtsiteration als quasi-zufälligen Prozess behandeln. In dieser Arbeit fassen wir die konzeptionelle Architektur des \emph{Euler-Adic-Borel-Cantor (EABC) / Bamberger Modells} zusammen. Durch die Einbettung des 2-adischen Bewertungselements in eine nicht-assoziative Quaternionen-Primzahlen-Geometrie und die Konstruktion expliziter, unimodularer Kettenbruch-Approximationszertifikate für $\beta = \log_2(9/8)$ leiten wir das Prinzip einer \emph{Vergangenheit-Zukunft-Trennung (Past--Future Separation)} her. Dieses erzwingt einen unlösbaren Widerspruch zwischen den exponentiellen Defekt-Unterrandwachstumsschranken (L3) und den linearen Majoranten fehlerfreier Null-Lift-Trajektorien (L4), wodurch unbegrenztes Wachstum global ausgeschlossen wird.
	\end{abstract}
	
	\tableofcontents % Korrigiert von \tableofdings
	\newpage
	
	\section{Einleitung und historischer Kontext}
	
	Die Collatz-Abbildung $T: \mathbb{N} \to \mathbb{N}$, definiert durch
	\begin{equation}
		T(n) = \begin{cases}
			\frac{n}{2} & \text{falls } n \text{ gerade,} \\
			\frac{3n+1}{2} & \text{falls } n \text{ ungerade,}
		\end{cases}
	\end{equation}
	präsentiert eine arithmetische Struktur von extremer Komplexität bei minimalistischer Definition. Obwohl numerische Verifizierungen die Abwesenheit nicht-trivialer Zyklen oder divergenter Orbits bis zu astronomischen Schranken ($>2\cdot 10^{20}$) belegen, blieb ein geschlossener analytischer Beweis fragmentiert. Die historischen Angriffsvektoren lassen sich in drei Epochen gliedern:
	\begin{enumerate}
		\item \textbf{Stochastische und empirische Analysen (Terras, Lagarias):} Die Untersuchung der erwarteten logarithmischen Drift zeigt zwar, dass Orbits im Mittel schrumpfen, kann jedoch seltene, maßtheoretisch verschwindende Divergenzpfade nicht absolut ausschließen.
		\item \textbf{Dyadische und $p$-adische Erweiterungen (Mahler, Everett):} Die Fortsetzung auf den Ring der $2$-adischen Zahlen $\mathbb{Z}_2$ enthüllte eine stetige Struktur, verbarg jedoch die feine algebraische Kopplung zwischen der multiplikativen Natur von $3$ und der metrischen Topologie von $2$.
		\item \textbf{Algebraisch-geometrische Starrheit (Das EABC-Modell):} Der hier verfolgte Ansatz bricht mit der isolierten Betrachtung einzelner Bahnen. Der globale Trajektorienraum wird als eine blätterteigartige Mannigfaltigkeit interpretiert, die von Quaternionen-Primzahlen, $E_8$-Gitterresonanzen und arithmetischen Kettenbrüchen regiert wird.
	\end{enumerate}
	
	---
	
	\section{Der arithmetische Kern: Irrationalität von $\beta$ und Kettenbruch-Zertifikate}
	
	Um die asymmetrische Bilanz aus Multiplikationen und Divisionen über beliebige Orbit-Horizonte zu kontrollieren, definieren wir die fundamentale Konstante:
	\begin{equation}
		\beta = \log_2 \left(\frac{9}{8}\right) = 2 \frac{\log 3}{\log 2} - 3.
	\end{equation}
	
	\begin{theorem}[Irrationalität von $\beta$]
		Die Zahl $\beta$ ist irrational.
	\end{theorem}
	\begin{proof}
		Wäre $\beta$ rational, so gäbe es ganze Zahlen $a, b > 0$ mit $\beta = a/b$, woraus eine nicht-triviale Relation $3^n = 2^m$ für positive Exponenten $n, m$ folgen würde. Dies widerspricht dem Satz von der eindeutigen Primfaktorzerlegung.
	\end{proof}
	
	Da $\beta$ irrational ist, bricht der reguläre Kettenbruch $\text{GenContFract.of } \beta$ niemals ab. Mithilfe der entsprechenden Mathlib-Infrastruktur konstruieren wir eine unbeschränkte Familie von Approximationszertifikaten.
	
	\begin{definition}[Positive-Side Approximation Certificate]
		Für jedes gerade $n = 2(k+1)$ liefern die Konvergenten $p_n/q_n$ und $p_{n+1}/q_{n+1}$ ein Zertifikat $(p, p', q, q', \varepsilon, \varepsilon')$ mit:
		\begin{enumerate}
			\item $0 < q < q'$ (strikte Nennerzunahme via Fibonacci-Untergrenzen),
			\item $p'q - pq' = 1$ (unimodulare Determinantenorientierung),
			\item $q\beta = p + \varepsilon$ mit $0 < \varepsilon < 1$,
			\item $q'\beta = p' - \varepsilon'$ mit $0 < \varepsilon' < \varepsilon$.
		\end{enumerate}
	\end{definition}
	
	Die Unbeschränktheit dieser Nennerfolge $q_k$ folgt direkt aus der Fibonacci-Untergrenze $n+1 \le \text{Fib}(n+1) \le q_n$ ($n \ge 5$) und der Wahl $k = M+2$.
	
	---
	
	\section{Die L3–L4-Architektur: Das Prinzip der Vergangenheit-Zukunft-Trennung}
	
	Das Kernargument verbindet die Kettenbruch-Zertifikate mit der dynamischen Wortstruktur der Collatz-Wand.
	
	\subsection{Der exponentielle Unterrand (L3)}
	Ein Zertifikat erzwingt im mechanischen Wort einen strukturellen Defekt bei einer präzisen Tiefe $m = q' - K - 1$. Dieser Defekt erzeugt eine Defektwand der Länge $L_m \ge 3m$, wodurch der Abstand zwischen gestörter und ungestörter Bahn nach $q$ Schritten exponentiell anwächst:
	\begin{equation}
		\|P_{K+q} - P_K\| \ge 2^{3m + 2}.
	\end{equation}
	
	\subsection{Die lineare Majorante (L4)}
	Unter der Null-Lift-Hypothese ($D(K) = 0$) treten im Intervall $[K, K+q)$ keine Defekte auf. Die Einzelschritte bleiben beschränkt ($\|P_{j+1} - P_j\| \le 1$), was über die Dreiecksungleichung zu einer streng linearen Schranke führt:
	\begin{equation}
		\|P_{K+q} - P_K\| \le \|P_K\| + q.
	\end{equation}
	
	\subsection{Der Widerspruch}
	Da die Kettenbruch-Geometrie die Schranke $q \le K + m$ erzwingt, hängen beide Seiten exakt von derselben Wachstumsvariablen $m$ ab. Für hinreichend große Defekttiefen dominiert die Exponentialfunktion jede affine Majorante:
	\begin{equation}
		\|P_K\| + K + m < 2^{3m + 2}.
	\end{equation}
	
	\begin{theorem}[Past--Future Separation]
		Für jede Shift-Tiefe $K \in \mathbb{N}$ gilt für den globalen Lift-Parameter:
		\begin{equation}
			D(K) \neq 0.
		\end{equation}
	\end{theorem}
	\begin{proof}
		Nimmt man an, dass $D(K) = 0$ gilt, so impliziert die Null-Lift-Propagation $D(K+q) = 0$ im gesamten Intervall. Die Wahl eines Zertifikats mit hinreichend großer Defekttiefe $m$ führt jedoch zum unmittelbaren Widerspruch zwischen L3 und L4. Folglich ist die Annahme $D(K) = 0$ unhaltbar.
	\end{proof}
	
	---
	
	\section{Formaler Status und offene Schnittstellen}
	
	Die formale Lean-4-Implementierung gliedert sich in folgende Module:
	\begin{itemize}
		\item \textbf{[M] Geschlossen:} Irrationalität von $\beta$, Nichttermination, unimodulare Determinanten, Residuen-Vorzeichen und unbeschränkte Zertifikatsfamilie.
		\item \textbf{[L4-closed] Geschlossen:} Intervallmechanik, Defekttiefen-Reduktion und lineare Akkumulation.
		\item \textbf{[L4-open-N1] bis [L4-open-N3] \& [L3-open]:} Die exakte rationale $D$-Rekurrenz, die Quantisierungsbrücke $D_K = 0 \implies \ell_K = 0$, die Einzelschritt-Kontrolle von $P$ sowie der exponentielle L3-Unterrand.
	\end{itemize}
	
	---
	
	\section{Fazit}
	
	\textbf{Die unerbittliche Wachstumsstarre der unimodularen Kettenbruch-Approximationen von $\beta$ schließt das Fortbestehen eines fehlerfreien Null-Lift-Zustands im dynamischen System kategorisch aus, wodurch die Existenz unbeschränkter Collatz-Trajektorien algebraisch unmöglich wird.}
	
\end{document}